\documentclass[11pt]{article}
\usepackage[margin=0.85in]{geometry}
\usepackage[T1]{fontenc}
\usepackage{lmodern,amsmath,graphicx,booktabs,longtable,xcolor,hyperref,microtype}
\hypersetup{colorlinks=true,linkcolor=blue,urlcolor=blue,citecolor=blue,
pdftitle={River-avulsion fingerprints from published image stacks and channel masks},
pdfauthor={Independent computational reproduction}}
\DisableLigatures{encoding=*,family=tt*}
\setlength{\parskip}{5pt}
\setlength{\parindent}{0pt}
\newcommand{\SR}{S_{\mathrm R}}
\newcommand{\figurefile}[3]{\IfFileExists{../figures/#1.pdf}{\begin{figure}[p]\centering
\includegraphics[width=\linewidth,height=.80\textheight,keepaspectratio]{../figures/#1.pdf}
\caption{#2}\label{#3}\end{figure}}{}}
\title{River-avulsion fingerprints from published image stacks and channel masks}
\author{Independent computational reproduction of Valenza et al. (2020)}
\date{October 2026}
\begin{document}
\maketitle
\begin{abstract}
River avulsions disturb different amounts of floodplain while establishing a
new flow path. Valenza et al. (2020) quantify this disturbance using annual
satellite observations and manually delineated channel masks. Here we test
the downstream calculation using the published processed images and masks,
with an independent C++ implementation and original plotting scripts.
This experiment accepts the authors' preprocessing and delineation as input;
it tests the subsequent spectral filtering, temporal accumulation and area
ratio calculation. Published measurements of channel geometry and morphology
are excluded from the analysis inputs. We calculate defined ratios for 63 events, spanning 0.03879--463.9, compared with the published range of 0.4--42.4. Thus the primary calculation does not reproduce the reported range. Explicit parent-window alternatives materially change the result. Independent geometry covers only four illustrative examples, so the population relationships remain untested.

The panel ledger and execution record identify the limits of this result. A successful software build
does not by itself establish reproduction of the scientific results.
\end{abstract}

\section{Question and scope}
The observable is the area of floodplain activity associated with channel
relocation. Its normalization by the area of the former channel defines a
dimensionless avulsion style ratio. A large ratio indicates a large disturbed
area relative to that channel; it does not by itself establish the mechanism
responsible for the disturbance.

The previous exercise began with independently acquired Landsat observations.
This exercise begins with the processed image stacks and parent/avulsion masks
deposited by Valenza and colleagues. The scientific specifications remain the
paper and its written supplementary material. No author MATLAB, Earth Engine,
analysis or plotting script is opened or executed. No published plot is
digitized, and no published figure is substituted for a generated result.

The user also permitted interpretation metadata: site identities and
coordinates, band/year mappings, selected bands and event-onset indices.
Published widths, slopes, distances, morphology classes, fingerprint areas,
style ratios and signature choices remain excluded from analysis inputs.
The metadata extractor reads only the authorized spreadsheet columns and
retains their source-row numbers. It does not retain the source workbook.

Independence is limited to the downstream implementation. The accepted masks
contain the authors' observational judgments about event extent and unrelated
water bodies. The processed images contain their cloud screening, annual
compositing and tasseled-cap transformation. Those operations are not newly
validated by matching a downstream fingerprint. Published numerical values
were seen during the earlier feasibility discussion and while extracting the
supplement as text. They were not copied into analysis inputs or used to tune
thresholds, windows or masks. This is not a blinded experiment.

\section{Data and mathematical specification}
The institutional archive supplies annual brightness, greenness and wetness
GeoTIFF stacks, together with region-of-interest masks. The archive period is
1986--2017. The paper describes Landsat observations with nominal 30 m sampling.
We preserve each supplied grid rather than silently imposing a different one.
The figure captions distinguish these accepted observations from quantities
calculated by the new implementation.

\subsection{Annual classification}
Let $B_y(x)$, $G_y(x)$ and $W_y(x)$ be the three transformed components at
pixel $x$ in annual layer $y$. We retain pixels with valid finite values in
all three components. File-defined nodata and validity masks are respected;
a finite zero is not automatically treated as missing. Across the retained
pixels in each complete image domain, calculate
\begin{align}
 I_{B,y}(x)&=\mathbf{1}\{B_y(x)>Q_{0.90}(B_y)\}
                  \mathbf{1}\{G_y(x)<Q_{0.50}(G_y)\},\\
 I_{W,y}(x)&=\mathbf{1}\{W_y(x)>Q_{0.75}(W_y)\}
                  \mathbf{1}\{G_y(x)<Q_{0.50}(G_y)\}.
\end{align}
Here $Q_p$ is the sample percentile. These rules identify unusually bright
or wet surfaces with low greenness. Both inequalities are strict. A constant
scene therefore contributes no detections.

We use linear interpolation at index $(n-1)p$ in a zero-indexed sorted sample
of $n$ valid values, retaining the convention used in the first exercise.
The written description does not specify a percentile interpolation rule.
The thresholds are recalculated from the supplied annual images and are
saved with valid-pixel and detection counts. They are not fitted to published
areas or ratios.

\subsection{Temporal accumulation and area}
For a signature $s$ and a channel region $r$, calculate
\begin{equation}
 F_{r,s}(x)=M_{r,s}(x)\sum_{y\in\mathcal{Y}_r} I_{s,y}(x),\qquad
 N_{r,s}=\sum_x\mathbf{1}\{F_{r,s}(x)>0\}.
\end{equation}
The accepted binary mask $M$ restricts activity to the indicated channel
region. A pixel detected in several years contributes once to the area
count, while its cumulative raster value records the number of detections.
We also save the number of valid observation years at each masked pixel.
Thus missing observations remain distinguishable from observed absence.

For a complete event, the style ratio is
\begin{equation}
 S_{\mathrm R,s}=\frac{N_{\mathrm{avulsion},s}}{N_{\mathrm{parent},s}}.
\end{equation}
A zero parent count gives an undefined ratio. Where both signatures have
adequate input coverage, the selected signature minimizes the avulsion
fingerprint area. This differs from minimizing the ratio when the parent
counts differ. Exact ties select brightness, consistent with the first
exercise. Both signature-specific counts remain available for inspection.

Pixel counts follow the published definition. The determinant of the image
transform is recorded as a native-coordinate pixel area. An angular pixel
area is not interpreted as square metres. Multiple image parts require
explicit accounting before their counts can be combined; overlap and missing
parts must not silently inflate or truncate an event.

\subsection{Time-window interpretation}
The authorized metadata provide a list of selected native bands and the
position of the onset of avulsion activity. The onset index refers to the
selected list, as shown by cases whose onset index does not occur in the
native-band list. With layers ordered from newest to oldest, positions up
to and including the onset index form the avulsion window; older positions
form the parent window. This interpretation is inferred from the written
metadata, without consulting the authors' program.

Some year labels disagree with the selected bands. These inconsistencies
are preserved in an audit table. Native band indices define the numerical
calculation; unresolved calendar labels are not silently replaced. Multiple
onset indices require separate timing scenarios. The exact band assignments
and any additional decisions are saved with the configuration.

\IfFileExists{metadata_notes.tex}{\begin{longtable}{p{.12\linewidth}p{.80\linewidth}}
\toprule Image part & Metadata issue\\\midrule\endhead
V14 & Year count differs from band count; intended year list unresolved.\\
B1\_b & Year count differs from band count; intended year list unresolved.\\
B14\_b & Multiple onset indices; evaluate both as separate timing scenarios.\\
H2 & Year count differs from band count; intended year list unresolved.\\
G4 & Invalid year range; preserve raw text and use supplied band indices. Year count differs from band count; intended year list unresolved.\\
G15 & Year label outside 1986--2017; label not used to change band selection.\\
\bottomrule\end{longtable}

Native image labels provide calendar names without changing the selected bands. They do not resolve which conflicting list the authors intended.
}{}

\subsection{Geometry and statistical relationships}
The paper's other observables are average channel width near the avulsion
site, along-channel distance to a physiographically identified mountain
front, parent-channel slope along a sufficiently long reach, and flow-thread
counts at multiple cross sections. The reported slope measurements usually
span 5--10 km and sometimes longer. A mask surrounding a channel is not a
measurement of these quantities.

The statistical implementation accepts independently documented geometry,
calculates normalized distance $D/w$, and fits a line in logarithmic space
using only $D/w>60$. Log-space ordinary least squares is an explicit
assumption because the paper does not fully specify its regression
convention. Group summaries use linearly interpolated quartiles and range
whiskers. Missing measurements remain absent. A spectral footprint alone
does not determine a channel's morphology class.

The four illustrative geometry observations are saved as explicit pixel-space
annotations. Width sampling integrates the wet-pixel intersections of three
short cross sections per event; a 1 m integration interval does not improve
the 30 m image resolution. A section must have valid observations and leave
the wet region at both ends. We assign a coarse morphology class only from
the independently recorded visible flow-thread counts. These observations
are single-date examples and may miss narrow threads.

We traced three example parent-channel or channel-belt centerlines before
calculating style ratios, sampled independently acquired Terrain Tiles
elevation at 30 m intervals, and estimated downstream slope by ordinary
least squares. We also save endpoint slopes. The illustrative logarithmic
comparison retains only profiles at least 5 km long with complete terrain
support and positive slope under both estimates. All attempted estimates,
including rejected ones, remain in the geometry table. This explicit
screen is a limited consistency check, not an elevation-error model.
Vegetation, DEM source differences, centerline placement and the short
sample prevent equating these estimates to the paper's slope measurements.

\section{Calculated results and panel assessment}
We calculated 77 image-part/timing jobs for 63 events.
63 events have defined style ratios under the stated primary assumptions.
The calculated range is 0.03879--463.9, from G23
to B10. Signature selection used the newly calculated avulsion
areas. It selected brightness for 57 events and
wetness for 6 events.

Every deposited event mask is named for only one spectral signature. We therefore
apply both filters to the same supplied geometry, retaining the mask's source name
only as provenance. This tests numerical signature selection on fixed geometry;
it cannot recreate a separately redrawn mask for the alternative signature.

The embedded image calendar labels disagree with 33 selected metadata
labels across 6 jobs. Native band selections
remain unchanged. The file \texttt{temporal\_labels.csv} records every assignment.
The unused-input ledger includes B5, which lacks a corresponding selected-window
record, and extra B10 tiles. No excluded measured result is used to fill these gaps.

The overlap audit identifies 15 site/signature/role combinations with
duplicate detected pixels. The primary event counts take the spatial union on
aligned native grids, so a physical pixel is counted once. A separate summed-part
calculation quantifies the effect of simply adding tile counts. Mosaic colors use
the maximum per-part detection count in overlaps; heterogeneous windows prevent
interpreting that color as one uniform observation period.

Across the 146 primary image-part/role windows, the fraction of
mask pixels observed at least once ranges from 91.59 to 100 percent;
22 windows have some never-observed pixels. This cumulative
coverage does not imply full coverage in every year. Missing pixels are not
filled, and the ratio is not corrected by an assumed missing-area multiplier.

We independently annotated 4 prespecified illustrative events
(B11, B12, V11 and V15). These are single-date, 30 m observations, not a new
63-event morphology survey. The elevation exercise produced 2 positive
slopes meeting the explicit reach-length and sign checks. Widths are wet-pixel
cross-section estimates with a native pixel scale of 30 m. No mountain-front
distance was measured. Consequently, the full distance relationship and the
published population distributions remain unreproduced.


\figurefile{fig1}{Reconstructed geographic site distributions for (a) the
northern Andes, (b) the central Andes, (c) the Himalaya and (d) New Guinea.
Coordinates are permitted source metadata; the terrain background comes
from the independent public terrain acquisition used in the first exercise.
B15 is included in this map. The map reproduces geographic content rather
than the publication's exact cartographic design.}{fig:locations}
\figurefile{fig1_locator}{Global study-site positions on geographic axes.
The continental outline is independent Natural Earth vector data. This
generated locator uses a simpler design than the published inset.}{fig:locator}
\figurefile{recomputed_ratios}{Style ratios calculated from the supplied
image stacks and masks by the new implementation. Each point is an event
with a defined ratio; sorting uses the computed ratio. Color denotes the
geographic prefix. The dashed line denotes equal avulsion and parent counts.
This additional diagnostic is not a substitute for the paper's geometry or
morphology relationships.}{fig:ratios}
\figurefile{fig2ab}{B12 full-event cumulative mosaic and illustrative detail from saved
independent calculations. Color measures the number of detected years.
Dashed contours, where present, outline the calculated parent fingerprint.
The input images and region masks are supplied data. Exact spatial and
temporal correspondence with the publication is not established: the full
mosaic is larger than the published central-part view, and the selected
detail is not the original validation window. Mosaic colors take the
maximum per-part count in overlaps.}{fig:b12}
\figurefile{fig2c}{Attempted annual counterparts for Figure 2c. The lower
row shows independently recomputed detections within the accepted avulsion
mask. The upper row records the availability of the separate high-resolution
validation observations. This independently chosen detail differs from the
published validation window, so these panels do not reproduce that
validation comparison. No published image is inserted as a replacement.}{fig:b12annual}
\figurefile{fig3}{Attempt to calculate the published style-ratio relationships.
Only independently documented geometry can define the horizontal coordinates.
An absence notice denotes an unreproduced relationship, not a null result.}{fig:relationships}
\figurefile{fig4}{Attempt to reconstruct style-ratio distributions by
independently measured morphology. Where data exist, boxes show quartiles
and whiskers show the full range. Missing classifications are not inferred
from the desired style ratio.}{fig:morphology}
\figurefile{fig5ce}{Independent fingerprint calculations for the two Sarare
events used in Figure 5. Color denotes detected years; dashed parent contours
are calculated from the corresponding observations. The accepted masks
remove the need for new manual event delineation in this exercise.}{fig:sarare}
\figurefile{fig5_context}{Generated context views for the Sarare events.
Tasseled-cap components are displayed as false color where available and
are explicitly distinguished from the paper's Google Earth imagery.}{fig:sararecontext}
\figurefile{fig5a}{Independent geographic locator for the two Sarare events.
The river course between the two deposited image domains has not been
traced, limiting correspondence to Figure 5a.}{fig:sararelocator}
\figurefile{geometry_observations}{Four independent illustrative observations.
Yellow lines mark the three width and flow-thread cross sections. Orange
lines mark the traced parent-channel or channel-belt centerline where one
was recorded. Backgrounds show the supplied TCT components with a
2--98 percent display stretch. Axes are native pixel coordinates; the
annotations and all derived measurements are saved separately.}{fig:geometry}
\figurefile{supp1}{Temporal context for B12. These views use the supplied
tasseled-cap components with a documented common display stretch. They are
not reconstructed conventional RGB photographs.}{fig:s1}
\figurefile{supp2}{Independent B11 fingerprint and detail, with the separate
high-resolution validation slot identified. Color records detected years.}{fig:s2}
\figurefile{supp3a}{Attempted slope distributions by independently measured
channel morphology. Missing observations are identified explicitly.}{fig:s3a}
\figurefile{supp3b}{Attempted normalized-distance distributions by
independently measured channel morphology.}{fig:s3b}
\clearpage

\section{Implementation, validation and reproducibility}
The C++17 implementation uses GDAL for raster input and output. Numerical
operations are implemented directly: percentile selection, strict spectral
classification, temporal accumulation, union counts, ratio selection and
statistical summaries. Python handles the authorized metadata, data-format
conversion, original plotting and report assembly. No author's program is
a build dependency.

Analytical checks cover constant scenes, missing values, singleton samples
and known quantiles. Separate raster tests exercise nodata, binary masks,
multiple years, temporal unions and rejection of shifted grids. Additional
checks use known multi-part totals and exact power-law relationships to
test signature selection and regression. These artificial fixtures are
software validation cases, never study data or manuscript results.

Real-data checks independently sum saved annual detections and compare the
result with the C++ cumulative rasters. They verify that detections are
binary, fingerprints do not exceed valid observation years, and saved pixel
counts agree with the rasters. These checks establish arithmetic consistency;
they cannot independently validate the supplied masks or observational
interpretation.

\IfFileExists{execution.tex}{The recorded full-rebuild state is \textbf{completed}. The recorded calculation took 630.0 s. The implementation passed 378 analytical quantile checks, 105 separate raster-integration checks, 11 summary/regression checks and 8 exact mosaic checks. The real outputs passed 6974 arithmetic and raster-consistency checks. The command-by-command execution record is \texttt{data/results/rebuild\_outcome.json}; logs and hashes are retained.
}{}

All project files, dependencies and caches are kept under a selected
\texttt{/data/data[n]} directory. The institutional raw archive is retrieved
through a documented download step and is not a substitute for saved
generated outputs. Input hashes, file inventories, configuration and commands
are retained. The README gives the exact complete rebuild procedure and
distinguishes a computation using cached inputs from a new download.

\section{Limits of the inference}
Accepting processed imagery and channel masks removes several decisions
that constrained the first exercise. Agreement after this input boundary
would test downstream numerical operations, while leaving the correctness
of the authors' preprocessing and manual masks outside the experiment.
Disagreement may arise from documented ambiguities in percentile conventions,
selected-band metadata, temporal splits, mask correspondence or multi-part
accounting. These alternatives should be assessed through explicit
sensitivity calculations, not by choosing the variant closest to a published
ratio.

The core physical interpretation additionally depends on channel geometry
and morphology. Recomputed fingerprints cannot establish those relationships
without the corresponding independent observations. The panel ledger
therefore separates calculation, partial correspondence and missing results.
\IfFileExists{comparison.tex}{\subsection{Quantitative comparison and sensitivity}

The paper reports a style-ratio range of 0.4--42.4 for 63 events. The present range is 0.03879--463.9. The lower and upper endpoints differ by -0.3612 and 421.5, respectively. These are distribution-endpoint comparisons, not matched-event errors. Matching the number of events does not establish matching event membership or measurements.

The published distance transition near $D/w=60$--$100$ and slope transition near $0.005$ are comparison targets only. They are not inputs to classification, masks, thresholds or geometry. The present independent geometry coverage cannot test either transition. The published individual-result and geometry tables are not imported.

\begin{center}\begin{tabular}{lrr}\toprule Calculation & Minimum $S_{\mathrm R}$ & Maximum $S_{\mathrm R}$\\\midrule

Primary spatial union & 0.03879 & 463.9\\

Sum of image-part counts & 0.03879 & 464\\

Alternative B14 onset & 0.03879 & 463.9\\

Parent in avulsion years & 0 & 30.86\\

Parent in all selected years & 0 & 24.85\\

\bottomrule\end{tabular}\end{center}

The temporal alternatives retain the same annual thresholds and avulsion windows; only the parent accumulation changes. They bracket specific interpretations of the written instruction to use the same time ranges where possible. They are not selected according to agreement with the paper. Every signature-specific result remains available.

The largest primary ratio illustrates the denominator sensitivity. For B10, the two calculated alternatives are B: 124780 avulsion pixels, 269 parent pixels and a ratio of 463.9; W: 520388 avulsion pixels, 25034 parent pixels and a ratio of 20.79. Minimizing avulsion area can choose a much larger ratio when that filter also detects far fewer parent-channel pixels. Both calculations use the same supplied mask geometry. This is an observed numerical consequence of the stated rule, not evidence that the larger ratio correctly measures floodplain disturbance.

The primary range fails to match the published range. Overlap accounting produces much smaller changes than parent-window selection. The experiment therefore establishes reproducibility of the stated downstream computation but does not establish reproduction of the paper's quantitative result. The missing independent geometry additionally prevents testing the proposed downstream transition.

For a limited matched-event check, Figure 5 explicitly prints $S_{\mathrm R}=0.5$ for V15 and $14.9$ for V11. These two numerical annotations were read only after the calculations, without extracting plotted coordinates or image areas. They enter this comparison section only; the numerical analysis never reads them.

For V15, the recomputed ratio is 2.77, a difference of 2.27 (454.0 percent) relative to the printed value.

For V11, the recomputed ratio is 28.66, a difference of 13.76 (92.3 percent) relative to the printed value.

Both calculations rank V11 above V15, but the magnitudes disagree. The brightness-selected V15 fingerprint is visibly sparse compared with the published depiction. The B12 full-event mosaic covers a larger domain than the published central-part view. Its independently chosen detail and annual views do not recover the original validation window; they demonstrate the calculation only. Display extents were refined after comparison, without changing any numerical input or geometry observation. The published images are absent from the generated figures and repository.
}{}

\appendix
\section{Figure-by-figure record}
\IfFileExists{panel_ledger.tex}{\small\begin{longtable}{p{.12\linewidth}p{.16\linewidth}p{.64\linewidth}}
\toprule Panel & Status & Evidence and limit\\\midrule\endhead
1a & reproduced & Permitted coordinates and independently acquired terrain/vector geography. Geographic content reproduced; cartographic design differs.\\
1b & reproduced & Permitted coordinates and independently acquired terrain/vector geography. Geographic content reproduced; cartographic design differs.\\
1c & reproduced & Permitted coordinates and independently acquired terrain/vector geography. Geographic content reproduced; cartographic design differs.\\
1d & reproduced & Permitted coordinates and independently acquired terrain/vector geography. Geographic content reproduced; cartographic design differs.\\
1 locator & reproduced & Permitted coordinates and independently acquired terrain/vector geography. Geographic content reproduced; cartographic design differs.\\
2a & partial & Independent cumulative raster from processed images and supplied masks. Timing, mask/signature convention and spatial correspondence limit equality claims.\\
2b & partial & Independent cumulative raster from processed images and supplied masks. Timing, mask/signature convention and spatial correspondence limit equality claims.\\
2c upper 2007 & not reproduced & Archive filename inventory checked. Separate WorldView validation scenes are absent.\\
2c upper 2010 & not reproduced & Archive filename inventory checked. Separate WorldView validation scenes are absent.\\
2c upper 2013 & not reproduced & Archive filename inventory checked. Separate WorldView validation scenes are absent.\\
2c upper 2017 & not reproduced & Archive filename inventory checked. Separate WorldView validation scenes are absent.\\
2c lower 2007 & partial & Annual detections recomputed from the selected native band. Independent display extent; exact reference-panel pixel agreement not measured.\\
2c lower 2010 & partial & Annual detections recomputed from the selected native band. Independent display extent; exact reference-panel pixel agreement not measured.\\
2c lower 2013 & partial & Annual detections recomputed from the selected native band. Independent display extent; exact reference-panel pixel agreement not measured.\\
2c lower 2017 & partial & Annual detections recomputed from the selected native band. Independent display extent; exact reference-panel pixel agreement not measured.\\
3a & not reproduced & Width observations attempted for four examples. No independently traced mountain-front distances; no distance regression.\\
3b & partial & Independent elevation profiles for illustrative examples. Insufficient coverage to reproduce the 63-event slope relationship.\\
4 & partial & Four independent visual classifications joined to computed ratios. Small prespecified subset; no population-wide three-class distribution.\\
5a & partial & Independent geographic locator for V11 and V15. Intervening river course and physiographic mountain-front boundary not traced.\\
5b & partial & Supplied TCT image rendered with a documented false-color stretch. Corresponding 2019 Google Earth imagery is absent; different observation date.\\
5c & partial & Independent cumulative raster from processed images and supplied masks. Timing, mask/signature convention and spatial correspondence limit equality claims.\\
5d & partial & Supplied TCT image rendered with a documented false-color stretch. Corresponding 2019 Google Earth imagery is absent; different observation date.\\
5e & partial & Independent cumulative raster from processed images and supplied masks. Timing, mask/signature convention and spatial correspondence limit equality claims.\\
S1a & partial & Requested-year processed components independently rendered. TCT false color cannot reconstruct the original conventional-color imagery.\\
S1b & partial & Requested-year processed components independently rendered. TCT false color cannot reconstruct the original conventional-color imagery.\\
S1c & partial & Requested-year processed components independently rendered. TCT false color cannot reconstruct the original conventional-color imagery.\\
S1d & partial & Requested-year processed components independently rendered. TCT false color cannot reconstruct the original conventional-color imagery.\\
S2a & partial & Independent cumulative raster from processed images and supplied masks. Timing, mask/signature convention and spatial correspondence limit equality claims.\\
S2b & partial & Independent cumulative raster from processed images and supplied masks. Timing, mask/signature convention and spatial correspondence limit equality claims.\\
S2c & not reproduced & Archive filename inventory checked. Separate high-resolution validation scene absent.\\
S3a & partial & Independent example slopes and visual morphology. No complete morphology-stratified population.\\
S3b & not reproduced & Geometric data requirements explicitly assessed. No mountain-front distances.\\
\bottomrule\end{longtable}\normalsize
}{}
\section{Computed event table}
\IfFileExists{event_table.tex}{\small\begin{longtable}{lrrrr}
\toprule Event & Filter & Avulsion pixels & Parent pixels & $S_{\mathrm R}$\\\midrule\endhead
B1 & B & 988498 & 50717 & 19.49\\
B10 & B & 124780 & 269 & 463.9\\
B11 & B & 649 & 347 & 1.87\\
B12 & B & 106667 & 2358 & 45.24\\
B13 & B & 297 & 346 & 0.8584\\
B14 & B & 392743 & 8212 & 47.83\\
B15 & B & 201 & 1143 & 0.1759\\
B2 & B & 156 & 290 & 0.5379\\
B3 & B & 82 & 451 & 0.1818\\
B4 & B & 198 & 13 & 15.23\\
B6 & B & 1036 & 891 & 1.163\\
B7 & B & 101 & 162 & 0.6235\\
B8 & W & 229 & 137 & 1.672\\
B9 & B & 33265 & 2964 & 11.22\\
G1 & B & 8688 & 5847 & 1.486\\
G10 & B & 22 & 63 & 0.3492\\
G11 & W & 260 & 193 & 1.347\\
G12 & B & 169 & 19 & 8.895\\
G13 & B & 669 & 18 & 37.17\\
G14 & B & 140 & 297 & 0.4714\\
G15 & B & 127 & 92 & 1.38\\
G16 & W & 5 & 10 & 0.5\\
G17 & B & 35 & 32 & 1.094\\
G18 & B & 57 & 97 & 0.5876\\
G19 & B & 375 & 180 & 2.083\\
G2 & B & 1124 & 11 & 102.2\\
G20 & B & 501 & 555 & 0.9027\\
G21 & B & 411 & 563 & 0.73\\
G22 & B & 377 & 592 & 0.6368\\
G23 & B & 9 & 232 & 0.03879\\
G24 & B & 391 & 480 & 0.8146\\
G25\_1 & B & 95 & 12 & 7.917\\
G25\_2 & B & 87 & 118 & 0.7373\\
G26 & W & 63 & 77 & 0.8182\\
G3 & B & 99 & 159 & 0.6226\\
G4 & B & 591 & 313 & 1.888\\
G5 & B & 1171 & 171 & 6.848\\
G6 & B & 365 & 174 & 2.098\\
G7 & B & 96 & 411 & 0.2336\\
G8 & B & 3190 & 38 & 83.95\\
G9 & B & 132 & 119 & 1.109\\
H1 & B & 21764 & 37876 & 0.5746\\
H2 & B & 3634 & 3533 & 1.029\\
H3 & W & 50 & 63 & 0.7937\\
H4 & W & 115 & 188 & 0.6117\\
H5 & B & 33145 & 36400 & 0.9106\\
H6 & B & 12000 & 1942 & 6.179\\
V1 & B & 270 & 170 & 1.588\\
V10 & B & 44 & 10 & 4.4\\
V11 & B & 14558 & 508 & 28.66\\
V12\_1 & B & 87 & 114 & 0.7632\\
V12\_2 & B & 376 & 7 & 53.71\\
V12\_3 & B & 22 & 306 & 0.0719\\
V13 & B & 59 & 1101 & 0.05359\\
V14 & B & 43914 & 155 & 283.3\\
V15 & B & 482 & 174 & 2.77\\
V4 & B & 95 & 30 & 3.167\\
V5\_1 & B & 138 & 295 & 0.4678\\
V5\_2 & B & 22 & 2 & 11\\
V6 & B & 1931 & 2395 & 0.8063\\
V7 & B & 2795 & 139 & 20.11\\
V8 & B & 579 & 573 & 1.01\\
V9 & B & 2227 & 2944 & 0.7565\\
\bottomrule\end{longtable}\normalsize
}{}
\section{Independent illustrative geometry}
\IfFileExists{geometry_table.tex}{\small\begin{tabular}{lrrrrr}
\toprule Event & Width (m) & Class & Reach (km) & OLS slope & Endpoint slope\\\midrule
B11 & unmeasured & 3 & 0 & undefined & undefined\\
B12 & 128.3 & 3 & 4.69 & -0.0004218 & 0.0004274\\
V11 & unmeasured & 3 & 7.686 & 0.001048 & 0.0007813\\
V15 & unmeasured & 1 & 7.983 & 0.007584 & 0.007393\\
\bottomrule\end{tabular}\normalsize\par
Classes are independent coarse observations: 1, multiple visible threads at all sections; 2, mixed single/multiple sections; 3, one visible thread at all sections. Width estimates require all three sections; missing values are retained. OLS and endpoint columns include unsuccessful slope attempts.
}{}

\begin{thebibliography}{9}
\bibitem{valenza} Valenza, J. M., Edmonds, D. A., Hwang, T. and Roy, S. (2020).
Downstream changes in river avulsion style are related to channel morphology.
\emph{Nature Communications} 11, 2116.
\url{https://doi.org/10.1038/s41467-020-15859-9}.
\bibitem{archive} Valenza, J., Edmonds, D., Hwang, T. and Roy, S. (2020).
Downstream changes in river avulsion style are related to channel morphology:
Data and Scripts. Indiana University DataCORE.
\url{https://hdl.handle.net/2022/25284};
\url{https://datacore.iu.edu/concern/data_sets/4x51hk339}.
Only the processed imagery, masks and permitted interpretation metadata
are inputs to this experiment.
\bibitem{terrain} Terrain Tiles public elevation archive. Provider
attribution and source metadata are recorded in
\texttt{data/provenance/terrain.json}.
\url{https://registry.opendata.aws/terrain-tiles/}.
\bibitem{land} Natural Earth, 1:110 million land polygons.
\url{https://www.naturalearthdata.com/downloads/110m-physical-vectors/110m-land/}.
Downloaded file identity is recorded in \texttt{data/provenance/land.json}.
\bibitem{first} First independent reproduction, source and compiled report:
\url{https://github.com/RianMontagh/valenza-2020-independent-reproduction}.
\end{thebibliography}
\end{document}
